\documentclass[11pt,a4paper]{article}

\usepackage[margin=2.5cm]{geometry}
\usepackage{amsmath,amssymb,amsthm}
\usepackage{bm}
\usepackage{booktabs}
\usepackage{algorithm}
\usepackage{algpseudocode}
\usepackage[round,authoryear]{natbib}
\usepackage[colorlinks=true,linkcolor=blue,citecolor=blue,urlcolor=blue]{hyperref}

\newtheorem{proposition}{Proposition}
\theoremstyle{definition}
\newtheorem{definition}{Definition}
\newtheorem{example}{Example}
\theoremstyle{remark}
\newtheorem*{remark}{Remark}

\newcommand{\R}{\mathbb{R}}
\newcommand{\T}{\mathcal{T}}
\newcommand{\ind}{\mathbf{1}}
\DeclareMathOperator{\Dir}{Dirichlet}
\DeclareMathOperator{\Gam}{Gamma}
\DeclareMathOperator{\Unif}{Uniform}
\DeclareMathOperator{\Expo}{Exponential}

\title{Sampling associativity fractions in \texttt{cntmosaic}:\\
a slice-within-Gibbs sampler on contingency tables}
\author{\texttt{cntmosaic} development notes}
\date{\today}

\begin{document}
\maketitle

\begin{abstract}
\noindent
This note explains how the function \texttt{sample\_eta} in \texttt{cntmosaic.predict} turns contact
counts that are known only by participant group into counts broken down by
both participant group and contact group. The method rests on one observation:
the two rules every valid answer must obey, \emph{reciprocity} and
\emph{sum-to-one}, turn each pair of age cells into a small table with known
row and column totals. We put a prior on such tables and sample from it with
a Markov chain Monte Carlo (MCMC) method that combines the classic
``swap'' moves for tables with fixed totals and the slice sampler of
\citet{neal2003slice}. Every step is explained from first principles.
\end{abstract}

\tableofcontents

%==============================================================================
\section{The problem}
\label{sec:problem}
%==============================================================================

\subsection{What we know and what we want}

People are split into $K$ groups, called \emph{strata} (for example
$K=2$: male and female), and into $A$ age groups. We write
$k,\ell \in \{1,\dots,K\}$ for strata and $a,b \in \{1,\dots,A\}$ for ages.
Throughout, $a$ is the age of the person \emph{reporting} the contact (the
participant) and $b$ the age of the person \emph{contacted}.

A partially stratified model, fitted to survey data where only the
participant's stratum is recorded, tells us
\begin{equation}
  z^{k}_{ab} = \text{total contacts from participants in stratum } k
  \text{ aged } a \text{ to people aged } b \text{ (any stratum)}.
\end{equation}
We want the finer quantity
\begin{equation}
  Z^{k\ell}_{ab} = \text{contacts from stratum } k \text{, age } a
  \text{, to stratum } \ell \text{, age } b,
\end{equation}
or equivalently the \emph{associativity fractions}
\begin{equation}
  \eta^{k\ell}_{ab} = \frac{Z^{k\ell}_{ab}}{z^{k}_{ab}},
\end{equation}
the share of stratum-$k$ participants' contacts (at ages $a \to b$) that go
to stratum $\ell$. In practice everything below is done separately for
each posterior draw of $z$; we drop the draw index to keep notation light.

\subsection{The rules a valid answer must obey}

\begin{enumerate}
  \item[(C1)] \textbf{Non-negativity:} $Z^{k\ell}_{ab} \ge 0$.
  \item[(C2)] \textbf{Sum-to-one:} the fine counts must add back up to what we
        know,
        \begin{equation}
          \sum_{\ell=1}^{K} Z^{k\ell}_{ab} = z^{k}_{ab}
          \quad\Longleftrightarrow\quad
          \sum_{\ell=1}^{K} \eta^{k\ell}_{ab} = 1 .
        \end{equation}
  \item[(C3)] \textbf{Reciprocity:} every contact has two ends. A contact
        from (stratum $k$, age $a$) to (stratum $\ell$, age $b$) is the same
        event as a contact from ($\ell$, $b$) to ($k$, $a$), so
        \begin{equation}
          Z^{k\ell}_{ab} = Z^{\ell k}_{ba}.
        \end{equation}
        Reciprocity is a standard requirement for contact matrices
        \citep{wallinga2006using,arregui2018projecting}.
\end{enumerate}

Summing (C3) over $k$ and $\ell$ and using (C2) shows that the known totals
must themselves be reciprocal:
\begin{equation}
  \label{eq:consistency}
  \sum_{k} z^{k}_{ab} = \sum_{k} z^{k}_{ba}.
\end{equation}
If \eqref{eq:consistency} fails, no answer satisfies (C1)--(C3). The
partially stratified models in \texttt{cntmosaic} satisfy it; the code warns
and rescales if the input is slightly off.

%==============================================================================
\section{The key observation: one small table per age pair}
\label{sec:tables}
%==============================================================================

Rules (C2) and (C3) only ever link a cell $(a,b)$ with its \emph{mirror}
cell $(b,a)$. So we can treat each pair of mirror cells on its own.

\subsection{Off-diagonal pairs ($a>b$)}

Fix $a > b$ and collect the unknown counts at $(a,b)$ in a $K\times K$
table $X$:
\begin{equation}
  X_{k\ell} = Z^{k\ell}_{ab}.
\end{equation}
\begin{itemize}
  \item By (C2) at $(a,b)$, the \textbf{row sums} of $X$ are known:
        $\sum_{\ell} X_{k\ell} = z^{k}_{ab} =: r_k$.
  \item By (C3), $X_{k\ell} = Z^{\ell k}_{ba}$, so the counts at the mirror
        cell are just the \textbf{transpose} $X^{\top}$.
  \item By (C2) at $(b,a)$, the rows of $X^{\top}$, i.e.\ the
        \textbf{column sums} of $X$, are known:
        $\sum_{k} X_{k\ell} = z^{\ell}_{ba} =: c_\ell$.
\end{itemize}
Condition \eqref{eq:consistency} says exactly that the row totals and the
column totals have the same grand total $N := \sum_k r_k = \sum_\ell c_\ell$.

\begin{definition}[Tables with fixed margins]
For non-negative vectors $\bm r, \bm c \in \R^{K}$ with the same total $N$,
\begin{equation}
  \T(\bm r,\bm c) = \Bigl\{ X \in \R^{K\times K} :
    X_{k\ell} \ge 0,\;
    \textstyle\sum_{\ell} X_{k\ell} = r_k,\;
    \sum_{k} X_{k\ell} = c_\ell \Bigr\}.
\end{equation}
This set is called a \emph{transportation polytope}.
\end{definition}

\begin{proposition}
\label{prop:equiv}
For a pair $a>b$, the counts at $(a,b)$ and $(b,a)$ satisfy (C1)--(C3) if and
only if $X \in \T(\bm r,\bm c)$ and the counts at $(b,a)$ equal $X^{\top}$.
\end{proposition}
\begin{proof}
This is the bullet list above read in both directions.
\end{proof}

So \emph{every} table in $\T(\bm r,\bm c)$ is a valid answer, and both
constraints hold exactly, by construction. Each entry automatically lies in
the classical Fr\'echet--Hoeffding bounds
\citep{hoeffding1940masstabinvariante,frechet1951tableaux}:
\begin{equation}
  \max(0,\; r_k + c_\ell - N) \;\le\; X_{k\ell} \;\le\; \min(r_k,\; c_\ell).
\end{equation}

\subsection{Diagonal cells ($a=b$)}

When $a=b$ the cell is its own mirror, so (C3) says $X_{k\ell} = X_{\ell k}$:
the table is \textbf{symmetric}, with row sums (and hence column sums)
$r_k = z^{k}_{aa}$.

\subsection{How much freedom is left?}

A $K\times K$ table with fixed row and column sums has $(K-1)^2$ free
entries; a symmetric one with fixed row sums has $K(K-1)/2$. For $K=2$ both
numbers are $1$: \emph{a single number decides everything}.

\begin{example}[$K=2$]
\label{ex:k2}
Let $x = X_{12}$ (stratum 1 to stratum 2). The margins fix the rest:
\begin{equation}
  X = \begin{pmatrix} r_1 - x & x \\ c_1 - r_1 + x & c_2 - x \end{pmatrix},
  \qquad
  x \in \bigl[\max(0,\, r_1 - c_1),\; \min(r_1,\, c_2)\bigr].
\end{equation}
Check: row 1 sums to $r_1$, column 1 to $c_1$, column 2 to $c_2$, and row 2
to $c_1 + c_2 - r_1 = r_2$ because the totals agree.
\end{example}

%==============================================================================
\section{The prior on tables}
\label{sec:prior}
%==============================================================================

\subsection{Starting point: one Dirichlet per source and cell}

Before imposing reciprocity, a natural prior treats the fractions of each
source stratum $k$ at each cell as a Dirichlet vector
\citep[see e.g.][]{ng2011dirichlet}:
\begin{equation}
  \bigl(\eta^{k1}_{ab},\dots,\eta^{kK}_{ab}\bigr) \sim
  \Dir\bigl(\alpha^{k1}_{ab},\dots,\alpha^{kK}_{ab}\bigr),
  \qquad
  \alpha^{k\ell}_{ab} = \alpha\, q^{\ell}_{b} + (\text{optional external terms}),
\end{equation}
where $q^{\ell}_{b}$ is the population share of stratum $\ell$ at age $b$ and
$\alpha>0$ is a precision: larger $\alpha$ pulls the fractions towards the
population shares. Recall that the Dirichlet density is
$\propto \prod_{\ell} \eta_\ell^{\alpha_\ell - 1}$ on the simplex.

\subsection{Imposing reciprocity}

At a pair $a>b$ the fractions at both cells are functions of $X$:
$\eta^{k\ell}_{ab} = X_{k\ell}/r_k$ and
$\eta^{k\ell}_{ba} = X_{\ell k}/c_k$. Multiplying the Dirichlet densities of
both cells and keeping only the tables in $\T(\bm r,\bm c)$ gives (the $r_k$
and $c_k$ are constants)
\begin{equation}
  \label{eq:prior}
  p(X) \;\propto\; \prod_{k,\ell}
  X_{k\ell}^{\,\beta_{k\ell} - 1},
  \qquad
  \beta_{k\ell} = \alpha^{k\ell}_{ab} + \alpha^{\ell k}_{ba} - 1,
  \qquad X \in \T(\bm r,\bm c).
\end{equation}
Each entry $X_{k\ell}$ appears in two places, $\eta^{k\ell}_{ab}$ and
$\eta^{\ell k}_{ba}$, so its exponent collects one term from each. For a
diagonal cell the same argument gives
$p(X) \propto \prod_{k} X_{kk}^{\alpha^{kk}_{aa}-1}
\prod_{k<\ell} X_{k\ell}^{\alpha^{k\ell}_{aa}+\alpha^{\ell k}_{aa}-2}$
on symmetric tables.

\begin{remark}[This is not a truncated Dirichlet]
A \emph{truncated} Dirichlet restricts the Dirichlet to a region of positive
volume (for example a box of bounds). Here we instead \emph{condition} on the
equality constraints (C3), a set of lower dimension. The resulting fractions
are in general neither Dirichlet nor truncated Dirichlet. Conditioning on a
set of probability zero depends on the coordinates used (the
Borel--Kolmogorov paradox; \citealp{kolmogorov1933grundbegriffe}); we use the
table entries, which are linear in the fractions, and the next proposition
shows this choice is natural.
\end{remark}

\begin{proposition}[A Gamma construction]
\label{prop:gamma}
Let $Y_{k\ell} \sim \Gam(\beta_{k\ell}, 1)$ be independent, with all
$\beta_{k\ell} > 0$. The conditional distribution of the table $Y$ given its
row sums $\bm r$ and column sums $\bm c$ has density \eqref{eq:prior}.
\end{proposition}
\begin{proof}
The joint density of $Y$ is
$\prod_{k\ell} y_{k\ell}^{\beta_{k\ell}-1} e^{-y_{k\ell}}/\Gamma(\beta_{k\ell})$.
On $\T(\bm r,\bm c)$ the sum $\sum_{k\ell} y_{k\ell} = N$ is fixed, so
$e^{-\sum y_{k\ell}}$ is a constant and only
$\prod y_{k\ell}^{\beta_{k\ell}-1}$ remains.
\end{proof}

This mirrors the familiar fact that a Dirichlet vector is a vector of
independent Gammas divided by their sum, i.e.\ conditioned on its total
\citep{ng2011dirichlet}; here we condition on all row \emph{and} column sums.
It is also the continuous analogue of the classical result that independent
Poisson counts conditioned on the margins of a table follow a (multivariate)
hypergeometric distribution, the setting of \citet{diaconis1998algebraic}.

\subsection{When is the prior proper?}

Near a face of $\T$ where $X_{k\ell} \to 0$, the density behaves like
$X_{k\ell}^{\beta_{k\ell}-1}$, which is integrable only if
$\beta_{k\ell} > 0$. With no external terms this reads
\begin{equation}
  \label{eq:proper}
  \beta_{k\ell} = \alpha\,\bigl(q^{\ell}_{b} + q^{k}_{a}\bigr) - 1 > 0
  \quad\Longleftrightarrow\quad
  \alpha > \frac{1}{q^{\ell}_{b} + q^{k}_{a}} .
\end{equation}
With two strata of roughly equal size the shares are near $1/2$, so we need
$\alpha$ somewhat above $1$; this is why the default is $\alpha = 2$.
The code checks \eqref{eq:proper} and raises an error that reports the
smallest admissible $\alpha$. (For $K=2$ one checks directly from
Example~\ref{ex:k2} that the condition is also necessary.)

%==============================================================================
\section{Sampling: Gibbs moves along margin-preserving directions}
\label{sec:gibbs}
%==============================================================================

We cannot draw from \eqref{eq:prior} directly: its normalising constant has
no closed form. We therefore use MCMC: build a Markov chain on $\T$ whose
long-run distribution is \eqref{eq:prior} \citep{robert2004monte}.

\subsection{Moves that keep the margins}

Write $E_{k\ell}$ for the table with a $1$ in position $(k,\ell)$ and zeros
elsewhere. For rows $k \ne k'$ and columns $\ell \ne \ell'$ define the
\emph{swap direction}
\begin{equation}
  M = E_{k\ell} + E_{k'\ell'} - E_{k\ell'} - E_{k'\ell}
  \qquad
  \begin{pmatrix} +1 & -1 \\ -1 & +1 \end{pmatrix}
  \text{ on rows } k,k' \text{ and columns } \ell,\ell'.
\end{equation}
Every row and column of $M$ sums to zero, so $X + dM$ has the same margins as
$X$ for every number $d$. These $2\times2$ moves are the classical tool for
exploring tables with fixed margins
\citep{diaconis1995rectangular,diaconis1998algebraic}. For symmetric
(diagonal) tables we use $M = E_{k\ell} + E_{\ell k} - E_{kk} - E_{\ell\ell}$,
which keeps both symmetry and the row sums.

Non-negativity limits how far we can move: $X + dM \in \T$ exactly when
\begin{equation}
  d \in [d_-, d_+], \qquad
  d_- = -\min(X_{k\ell},\, X_{k'\ell'}), \qquad
  d_+ = \min(X_{k\ell'},\, X_{k'\ell}).
\end{equation}

\subsection{The one-dimensional target}

Along the move, the density \eqref{eq:prior} becomes a function of $d$ only:
\begin{equation}
  \label{eq:onedim}
  g(d) \;\propto\;
  (X_{k\ell} + d)^{\beta_{k\ell}-1}\,
  (X_{k'\ell'} + d)^{\beta_{k'\ell'}-1}\,
  (X_{k\ell'} - d)^{\beta_{k\ell'}-1}\,
  (X_{k'\ell} - d)^{\beta_{k'\ell}-1},
  \qquad d \in [d_-, d_+].
\end{equation}
Drawing $d$ from $g$ and setting $X \leftarrow X + dM$ is a \emph{Gibbs
update along a direction} \citep{gelfand1990sampling}, closely related to
the hit-and-run sampler of \citet{smith1984efficient}. Each such update
leaves \eqref{eq:prior} invariant, and so does any fixed sequence of them
\citep[Section~2.4]{tierney1994markov}. One \emph{sweep} applies every swap
direction once; the swap directions span every direction in which a table in
$\T$ can move, so repeated sweeps can explore all of $\T$.

For $K=2$ there is exactly one direction, and $d$ is simply the change in
$x = X_{12}$ of Example~\ref{ex:k2}.

%==============================================================================
\section{The slice sampler}
\label{sec:slice}
%==============================================================================

\subsection{Why not sample $g$ directly?}

$g$ is a product of up to four powers of linear functions. It is not a
standard distribution, and neither its cumulative distribution function nor
its inverse has a closed form, so we cannot use the inverse-CDF method.
Common alternatives have drawbacks here:
\begin{itemize}
  \item \textbf{Random-walk Metropolis} needs a step size, and the right size
        varies enormously across the thousands of age pairs and posterior
        draws we update at once.
  \item \textbf{Rejection sampling} needs an envelope function that lies above
        $g$. When an exponent $\beta-1$ is between $-1$ and $0$, $g$ is
        infinite at an end of the interval (though still integrable), which
        makes envelopes awkward.
\end{itemize}
The slice sampler of \citet{neal2003slice} avoids both problems: it has no
step size to tune, copes with infinite spikes at the ends, needs $g$ only up
to a constant, and leaves $g$ exactly invariant.

\subsection{The idea}

Sampling from a density is the same as picking a point uniformly from the
region under its graph, $\{(d,u): 0 < u < g(d)\}$, and keeping the $d$
coordinate \citep{neal2003slice,damien1999gibbs}. The slice sampler does this
by alternating two easy steps from the current point $d_0$:
\begin{enumerate}
  \item \textbf{Pick a height.} Draw $u \sim \Unif(0, g(d_0))$.
  \item \textbf{Pick a point on the slice.} Draw the new $d$ uniformly from
        the \emph{slice} $S = \{d : g(d) > u\}$.
\end{enumerate}
Both steps are Gibbs updates for the uniform distribution under the graph,
so the pair $(d,u)$ keeps that distribution, and $d$ keeps distribution $g$.

For numerical safety we work with logarithms: step 1 is equivalent to
$y = \log g(d_0) - \varepsilon$ with $\varepsilon \sim \Expo(1)$, and the
slice is $S = \{d : \log g(d) > y\}$.

\subsection{Finding a point on the slice: shrinkage}

We usually cannot write $S$ down, but we can test whether a given $d$ is in
it. Neal's \emph{shrinkage} procedure \citep[Section~4.2]{neal2003slice}
starts from an interval $[L,R]$ that contains the whole slice and repeats:
\begin{enumerate}
  \item propose $d^\star \sim \Unif(L,R)$;
  \item if $d^\star \in S$, accept it and stop;
  \item otherwise shrink the interval towards the current point: if
        $d^\star < d_0$ set $L = d^\star$, else set $R = d^\star$.
\end{enumerate}
Because the current point $d_0$ is always in $S$ and always stays inside
$[L,R]$, the loop ends. Neal proves that this procedure leaves the uniform
distribution on $S$ invariant, even when $S$ is not a single interval
\citep[Sections~4.2--4.3]{neal2003slice}; this matters here because $g$ can
be U-shaped, giving a slice made of two pieces.

Our situation is the easiest one Neal describes: $g$ lives on the bounded
interval $[d_-, d_+]$, so we start with $[L,R] = [d_-, d_+]$, which certainly
contains $S$. No ``stepping-out'' phase is needed.

\subsection{The full algorithm}

\begin{algorithm}[ht]
\caption{One slice update of a table $X$ along direction $M$
(current step $d_0 = 0$)}
\label{alg:slice}
\begin{algorithmic}[1]
\State $d_- \gets -\min(X_{k\ell}, X_{k'\ell'})$,\;
       $d_+ \gets \min(X_{k\ell'}, X_{k'\ell})$
\If{$d_- = d_+$} \Return $X$ \Comment{no room to move}
\EndIf
\State $y \gets \log g(0) - \varepsilon$, \quad $\varepsilon \sim \Expo(1)$
\State $L \gets d_-$, \quad $R \gets d_+$
\Loop
  \State $d^\star \sim \Unif(L, R)$
  \If{$\log g(d^\star) > y$}
    \State \Return $X + d^\star M$
  \ElsIf{$d^\star < 0$}
    \State $L \gets d^\star$
  \Else
    \State $R \gets d^\star$
  \EndIf
\EndLoop
\end{algorithmic}
\end{algorithm}

\begin{algorithm}[ht]
\caption{\texttt{sample\_eta}: one table per age pair and posterior draw}
\label{alg:full}
\begin{algorithmic}[1]
\For{each posterior draw and each age pair $a \ge b$}
  \State $\bm r \gets (z^{k}_{ab})_k$,\; $\bm c \gets (z^{\ell}_{ba})_\ell$,\;
         $N \gets \sum_k r_k$
  \State $X \gets \bm r \bm c^{\top} / N$
         \Comment{independence table: strictly inside $\T$}
\EndFor
\For{sweep $= 1, \dots, n_{\text{sweeps}}$}
  \For{each swap direction $M$ (symmetric directions on the diagonal)}
    \State update every table with Algorithm~\ref{alg:slice}
  \EndFor
\EndFor
\State $\eta^{k\ell}_{ab} \gets X_{k\ell}/r_k$,\;
       $\eta^{k\ell}_{ba} \gets X_{\ell k}/c_k$
\end{algorithmic}
\end{algorithm}

%==============================================================================
\section{Implementation notes}
\label{sec:impl}
%==============================================================================

\begin{itemize}
  \item \textbf{Starting point.} The independence table $\bm r\bm c^{\top}/N$
        has the right margins and positive entries whenever the margins are
        positive, so it is a valid starting point. MCMC draws are only exact
        in the limit; after a finite number of sweeps they still depend
        slightly on where the chain started.
  \item \textbf{Number of sweeps.} The default is
        $n_{\text{sweeps}} = 10$. In a test cell with an exactly computable
        mean, $1$--$3$ sweeps were visibly biased towards the starting point,
        while $5$ or more sweeps matched the exact mean.
  \item \textbf{Vectorisation.} All tables (every posterior draw and age
        pair) run as independent chains, updated together with NumPy. Almost
        every chain accepts its first proposal (on average $1.04$ proposals
        per update), so after the first proposal the code only keeps working
        on the few chains that have not yet accepted. This gave a
        $6\times$ speed-up without changing the algorithm.
  \item \textbf{Zero margins.} If a source has no contacts at a cell
        ($r_k = 0$), its row is zero, every move involving it has
        $d_- = d_+ = 0$, and its fractions are reported as $0$.
\end{itemize}

%==============================================================================
\section{Checks}
\label{sec:checks}
%==============================================================================

The unit tests in \texttt{cntmosaic/predict/tests/test\_predict.py} check:
\begin{itemize}
  \item sum-to-one (C2) and reciprocity (C3) hold to machine precision
        (around $10^{-16}$), for $K=2$ and $K=3$;
  \item every fraction lies within its Fr\'echet--Hoeffding bounds;
  \item for $K=2$, the sampled mean of $x$ in Example~\ref{ex:k2} agrees,
        within four standard errors, with the exact mean of \eqref{eq:onedim}
        computed by numerical integration;
  \item an improper prior (condition \eqref{eq:proper} violated) raises an
        error, and non-reciprocal input totals trigger a warning.
\end{itemize}

%==============================================================================
\section{Summary and limitations}
\label{sec:summary}
%==============================================================================

\begin{itemize}
  \item Reciprocity and sum-to-one turn each age pair into a table with known
        row and column totals (Proposition~\ref{prop:equiv}); any such table is
        a valid answer.
  \item The prior is the Dirichlet prior of each cell, conditioned on
        reciprocity; equivalently, independent Gamma entries conditioned on
        the margins (Proposition~\ref{prop:gamma}). It is proper when
        condition \eqref{eq:proper} holds.
  \item We sample it with Gibbs updates along margin-preserving swap
        directions \citep{diaconis1998algebraic}, each drawn exactly with a
        shrinkage slice sampler \citep{neal2003slice}.
\end{itemize}

\noindent
\textbf{Limitations.}
(i) The draws are approximate for a finite number of sweeps.
(ii) The conditioned prior depends on the choice of coordinates (see the
Remark in Section~\ref{sec:prior}); we use the table entries.
(iii) We are not aware of a published source that uses exactly this
construction for contact matrices; the components above are standard, and
the closest analogue is the Diaconis--Sturmfels approach for discrete tables.

%==============================================================================
\begin{thebibliography}{99}
%==============================================================================

\bibitem[Arregui et al.(2018)]{arregui2018projecting}
Arregui, S., Aleta, A., Sanz, J. and Moreno, Y. (2018).
Projecting social contact matrices to different demographic structures.
\emph{PLoS Computational Biology}, 14(12), e1006638.

\bibitem[Damien et al.(1999)]{damien1999gibbs}
Damien, P., Wakefield, J. and Walker, S. (1999).
Gibbs sampling for Bayesian non-conjugate and hierarchical models by using
auxiliary variables.
\emph{Journal of the Royal Statistical Society: Series B}, 61(2), 331--344.

\bibitem[Diaconis and Gangolli(1995)]{diaconis1995rectangular}
Diaconis, P. and Gangolli, A. (1995).
Rectangular arrays with fixed margins.
In \emph{Discrete Probability and Algorithms}, IMA Volumes in Mathematics and
its Applications, vol.~72, 15--41. Springer.

\bibitem[Diaconis and Sturmfels(1998)]{diaconis1998algebraic}
Diaconis, P. and Sturmfels, B. (1998).
Algebraic algorithms for sampling from conditional distributions.
\emph{The Annals of Statistics}, 26(1), 363--397.

\bibitem[Fr\'echet(1951)]{frechet1951tableaux}
Fr\'echet, M. (1951).
Sur les tableaux de corr\'elation dont les marges sont donn\'ees.
\emph{Annales de l'Universit\'e de Lyon, Section A}, 14, 53--77.

\bibitem[Gelfand and Smith(1990)]{gelfand1990sampling}
Gelfand, A. E. and Smith, A. F. M. (1990).
Sampling-based approaches to calculating marginal densities.
\emph{Journal of the American Statistical Association}, 85(410), 398--409.

\bibitem[Hoeffding(1940)]{hoeffding1940masstabinvariante}
Hoeffding, W. (1940).
Masstabinvariante Korrelationstheorie.
\emph{Schriften des Mathematischen Instituts und des Instituts f\"ur
Angewandte Mathematik der Universit\"at Berlin}, 5, 181--233.

\bibitem[Kolmogorov(1933)]{kolmogorov1933grundbegriffe}
Kolmogorov, A. N. (1933).
\emph{Grundbegriffe der Wahrscheinlichkeitsrechnung}. Springer, Berlin.

\bibitem[Neal(2003)]{neal2003slice}
Neal, R. M. (2003).
Slice sampling.
\emph{The Annals of Statistics}, 31(3), 705--767.

\bibitem[Ng et al.(2011)]{ng2011dirichlet}
Ng, K. W., Tian, G.-L. and Tang, M.-L. (2011).
\emph{Dirichlet and Related Distributions: Theory, Methods and
Applications}. Wiley.

\bibitem[Robert and Casella(2004)]{robert2004monte}
Robert, C. P. and Casella, G. (2004).
\emph{Monte Carlo Statistical Methods}, 2nd edition. Springer.

\bibitem[Smith(1984)]{smith1984efficient}
Smith, R. L. (1984).
Efficient Monte Carlo procedures for generating points uniformly distributed
over bounded regions.
\emph{Operations Research}, 32(6), 1296--1308.

\bibitem[Tierney(1994)]{tierney1994markov}
Tierney, L. (1994).
Markov chains for exploring posterior distributions.
\emph{The Annals of Statistics}, 22(4), 1701--1762.

\bibitem[Wallinga et al.(2006)]{wallinga2006using}
Wallinga, J., Teunis, P. and Kretzschmar, M. (2006).
Using data on social contacts to estimate age-specific transmission
parameters for respiratory-spread infectious agents.
\emph{American Journal of Epidemiology}, 164(10), 936--944.

\end{thebibliography}

\end{document}
